The goal of this project was not to achieve a specific level of accuracy with the neural network implemented on the FPGA, but rather to successfully implement the network and obtain an inference time low enough to be unnoticeable to a human user. \\
This requirement was achieved using a complete software implementation while accelerating the floating-point arithmetic operations through dedicated hardware. \\
If this implementation had not been sufficient to achieve an adequately low inference time, several improvements could have been introduced.

\subsection{Possible upgrades}

The possible upgrades are various, but the most important metrics to improve are:

\begin{itemize}
    \setlength\itemsep{0pt}
    \setlength\parskip{0pt}
    \item \textbf{Inference Time}: reducing the inference time is important to compute the output faster.
    \item \textbf{Average Accuracy}: improving the accuracy of the network allows more characters to be recognized correctly on average.
\end{itemize}

\paragraph{Inference time}

In \autoref{sssec:nn_inf_time}, some optimizations and their impact on the inference time of the network are presented. \\
However, if even with these improvements the inference time is not sufficiently small, an alternative approach would be to describe the entire network in hardware using an HDL language and integrate it as a custom module of the CPU. \\
This method is more complex and less modular than a software-based approach, but it can achieve significantly higher speed. \\
If even this solution is not sufficient to meet the inference time requirements, dedicated external hardware accelerators should be used to improve multiply-and-accumulate operations (like GPUs). \\
Another possible approach is to quantize the weights and biases into integer values, simplifying multiplication and addition operations; however, this may lead to reduced accuracy.

\paragraph{Average Accuracy}

The network average accuracy with MNIST dataset, is approximately the \(85\%\), but this value is computed by using some letters from the dataset; the letters wrote on the screen are usually different in orientation and dimensions with the one on the dataset, so the resulting accuracy is less than the one computed (the \(85\%\) accuracy computed during training, see \autoref{sssec:nn_accuracy_tst}). \\
To increase the real accuracy some precaution has been taken, like the one described in the \textit{Note} at the end of \autoref{sssec:nn_accuracy_tst}, but the accuracy remains always lower than \(85\%\). \\
During the test has been tried different types of neural networks, with more hidden layers, or with less hidden layers but with more neurons, but the accuracy increased to the utmost at \(86\%\), the one used for the project is the best compromise of memory-footprint and accuracy of the network. \\
To rise the accuracy a different type of network must be used, like convolutional networks or combination of convolutional and fully-connected networks like LeNet-5 (Reference \cite{lecun1998}). \\
Another limitation is that the network performs poorly when recognizing distorted or scaled letters and numbers. \\
To reduce this problem, the network could be trained using an enlarged dataset generated by applying transformations such as distortion, scaling, and rotation to the original characters.

\vspace{3em}
\noindent
\textit{Note: To compute the real accuracy of the network, different letters would need to be written manually and compared with the expected results. \\
This process can be performed, but it is time-consuming and was not carried out in this project.}

\subsection{Implementation on Different FPGA Devices}

As can be observed in \autoref{sssec:hw_footprint_chpl}, the accelerated floating-point arithmetic project instantiates a significant amount of hardware resources. In particular, the most heavily utilized resources are the embedded 9-bit multipliers: the design instantiates \(16\) of them. \\
Another largely utilized resource is the on-chip memory, with a total usage of \(180,992\) memory bits. \\
These resources are commonly available on most modern FPGAs.

\begin{itemize}
    \setlength\itemsep{0pt}
    \setlength\parskip{0pt}
    \item The MAX10 FPGA (the one we used for this project) has \(288\) available embedded 9bit-multiplier, so their usage for this project is approximately the \(5.5\%\), and \(1,677,312\) memory bits, corresponding to a usage of approximately the \(10.8\%\)
    \item The CycloneII, an Altera older FPGA, has \(70\) available embedded 9bit-multiplier,  so their usage for this project is approximately the \(22.8\%\), and \(483,840\) memory bits, corresponding to a usage of approximately the \(37.4\%\)
\end{itemize}

\noindent
The relatively low resource utilization allows the proposed architecture to be implemented even on older FPGA devices. \\
However, timing analysis may still report violations at the target operating frequency due to technological differences and the lower performance of older FPGA architectures. \\
This results in a lower maximum operating frequency and, consequently, a longer inference time. \\
In some cases, especially in the presence of hold-time violations, the architecture may fail to operate correctly even when the clock frequency is reduced.
